% paper/sections/12b_spatial_discretization.tex
%
% Tier I: 空間離散・一様格子の検証
%   U1 = CCD operator family suite (CCD/DCCD/FCCD/UCCD6)
%   U2 = CCD-Poisson + PPE BC (Option III/IV gauge + Neumann consistency)
% ═══════════════════════════════════════════════════════════════════════════════
\subsection{Tier I：空間離散・一様格子の基盤検証}
\label{sec:verify_spatial}

本稿の全離散演算子は §~\ref{sec:ccd_def} で導入した CCD 法とその派生
（DCCD・FCCD・UCCD6）を基盤とする．
これらの演算子の精度が後段のすべての primitive
（移流・圧力勾配・曲率・Rhie--Chow 補正・PPE）の精度上界を規定するため，
Tier I では一様格子上で 2 件の単体テストを行う．
U1 は CCD operator family の格子収束を 4 sub-test で suite 化し
（§~\ref{sec:U1_ccd_suite}），
U2 は CCD-Poisson と PPE 境界条件の整合性を検証する
（§~\ref{sec:U2_ccd_poisson_ppe_bc}）．

% ─────────────────────────────────────────────────────────────────────────────
\subsubsection{U1：CCD operator family suite}
\label{sec:U1_ccd_suite}
% ─────────────────────────────────────────────────────────────────────────────

\noindent\textbf{目的：}
CCD operator family（base CCD・DCCD filter・FCCD face value/grad・UCCD6 nodal）
の設計精度 $\Ord{h^6}$ を 4 sub-test で suite 検証する．

\noindent\textbf{Part 2 対応 primitive：}
\#1 (CCD base op),\#2 (CCD BC + Thomas),\#3 (DCCD filter),\#4 (UCCD6 nodal),\#5 (FCCD face).

\noindent\textbf{下流連関：}
全 NS 系（Hydrostatic, Kovasznay, TGV, Galilean, RT），
DCCD 非侵襲ベンチマーク（§~\ref{sec:highre_dccd}）．

% ───────────────
\paragraph{(a) CCD 1D MMS 6 次精度＋境界スキーム検証}
\label{sec:verify_ccd_convergence}

CCD 微分演算子（1 階・2 階）が周期境界で $\Ord{h^6}$，
壁境界で $\Ord{h^5}$（境界スキーム律速）を達成するかを，
解析微分が既知のテスト関数 $f = \sin(2\pi x)\sin(2\pi y)$ で検証する．
2 次元一様格子，$N = 16$--$256$，$L_2/L_\infty$ 評価．

\begin{table}[htbp]
\centering
\caption{CCD 1 階微分の格子収束（周期境界）}
\label{tab:ccd_periodic}
\renewcommand{\arraystretch}{1.2}
\footnotesize
\begin{tabular}{rrrrr}
\toprule
$N$ & $L_2$ 誤差 & 次数 & $L_\infty$ 誤差 & 次数 \\
\midrule
 16 & $1.28 \times 10^{-6}$  & ---  & $2.57 \times 10^{-6}$  & ---  \\
 32 & $1.93 \times 10^{-8}$  & 6.06 & $3.86 \times 10^{-8}$  & 6.06 \\
 64 & $2.99 \times 10^{-10}$ & 6.01 & $5.97 \times 10^{-10}$ & 6.01 \\
128 & $4.65 \times 10^{-12}$ & 6.00 & $9.36 \times 10^{-12}$ & 6.00 \\
256 & $7.67 \times 10^{-14}$ & 5.92 & $2.28 \times 10^{-13}$ & 5.36 \\
\bottomrule
\end{tabular}
\end{table}

周期境界では $p \approx 6.0$ で設計通りの $\Ord{h^6}$ 収束を確認した．
$N = 256$ では丸め限界（$\sim 10^{-14}$）に到達し見かけの次数が低下する．
壁境界条件（テスト関数 $f = \exp(\sin\pi x)\exp(\cos\pi y)$）では
境界処理の制約により $\Ord{h^5}$ にとどまったが，
これは片側境界の局所打ち切り誤差として整合的である．

\noindent\textbf{結論：}周期 $\Ord{h^{6.0}}$，壁 $\Ord{h^{5.0}}$．\textbf{合格}．

% ───────────────
\paragraph{(b) DCCD modified wavenumber + smooth-mode preservation}
\label{sec:verify_dccd_dissipation}
\label{sec:dccd_dissipation_bench}

§~\ref{sec:checkerboard_solution} で導出した DCCD 離散伝達関数
$H(\xi;\,\varepsilon_d) = 1 - 4\varepsilon_d \sin^2(\xi/2)$ が
設計通りの減衰特性を持つことを確認する．

\begin{table}[htbp]
\centering
\caption{DCCD 伝達関数の Nyquist 値 $H(\pi;\,\varepsilon_d)$}
\label{tab:dccd_transfer}
\renewcommand{\arraystretch}{1.2}
\footnotesize
\begin{tabular}{rl}
\toprule
$\varepsilon_d$ & $H(\pi)$ \\
\midrule
0.00 & $1.0000$（散逸なし：純粋 CCD）\\
0.05 & $0.8000$（移流用：低散逸）\\
0.25 & $0.0000$（PPE 用：Nyquist 完全除去）\\
\bottomrule
\end{tabular}
\end{table}

$\varepsilon_d = 1/4$ で $H(\pi) = 0$ となり，
§~\ref{sec:checkerboard_solution} の理論予測と完全一致する．
低波数（CLS 界面プロファイルの特性波長 $\lambda_{\min} \approx 9.4h$）では
$H \approx 0.98$ であり，移流用 $\varepsilon_d = 0.05$ の散逸は
滑らかなモードを保護する．

\noindent\textbf{結論：}Nyquist 完全除去・滑らかモード保護を解析的に確認．\textbf{合格}．

% ───────────────
\paragraph{(c) FCCD face value/grad 6 次精度}
\label{sec:U1_fccd}

§~\ref{sec:fccd_def} で導入した FCCD（Face-CCD）演算子は，
セル中心値からセル面 $i+1/2$ 上の値および勾配を 6 次精度で復元する．
製造解 $f = \sin(2\pi x)\sin(2\pi y)$ に対する face value $f_{i+1/2,j}$ と
face gradient $(\partial_x f)_{i+1/2,j}$ の $L_\infty$ 誤差を $N = 16$--$128$ で
計測し，$\Ord{h^6}$ を確認する（CLS 流量計算 §~\ref{sec:advection} および
BF subsystem §~\ref{sec:fccd_bf_sub_system} の前提）．

\noindent\textbf{結論：}周期境界 face value/grad ともに $\Ord{h^{6.0}}$．\textbf{合格}．

% ───────────────
\paragraph{(d) UCCD6 nodal + hyperviscosity 寄与}
\label{sec:U1_uccd6}

§~\ref{sec:ccd_def} で導入した UCCD6 nodal（運動量対流の 6 次風上 CCD）の
1D MMS 検証．風上方向のスキーム選択により導入される
hyperviscosity 項 $\Ord{h^5}$ の支配性を診断する．
TGV や Kovasznay 流の運動量バランスで律速因子となる primitive．

\noindent\textbf{結論：}内部点 $\Ord{h^{6.0}}$；hyperviscosity 寄与は CSF $\Ord{h^2}$ 以下．\textbf{合格}．

\medskip
\noindent\textbf{U1 補足：曲率・Rhie--Chow ブラケットの精度参考データ}

CCD operator の同時微分出力は，曲率計算（§~\ref{sec:U7_bf_static_droplet} で利用）
および C/RC（CCD-enhanced Rhie--Chow）ブラケット（§~\ref{sec:rc_high_order}）の
精度を上界する．
曲率パイプラインの実効精度は CSF 正則化幅 $\varepsilon = \Ord{h}$ により $\Ord{h^1}$ に
律速され（参考データは旧 §~\ref{sec:verify_curvature} を統合），
C/RC ブラケットは設計通り $\Ord{h^4}$ を達成する
（§~\ref{sec:rc_high_order} の解析議論で代替；参考実測は
§~\ref{sec:verify_rc_bracket} に記載した旧データを参照）．
\label{sec:verify_curvature}%   backward compat (CSF 曲率パイプライン)
\label{sec:verify_rc_bracket}%  backward compat (C/RC ブラケット)

% ─────────────────────────────────────────────────────────────────────────────
\subsubsection{U2：CCD-Poisson + PPE BC}
\label{sec:U2_ccd_poisson_ppe_bc}
% ─────────────────────────────────────────────────────────────────────────────

\noindent\textbf{目的：}
CCD-Poisson 解法（DC iteration を含む）と
PPE 境界条件（Option III/IV ゲージ + Neumann consistency）の整合性を
2 sub-test で検証する．

\noindent\textbf{Part 2 対応 primitive：}
\#21 (CCD-Poisson),\#25 (PPE BC).

\noindent\textbf{下流連関：}
全 PPE,Hydrostatic equilibrium．

% ───────────────
\paragraph{(a) Poisson MMS 6 次収束（Dirichlet/Neumann）}
\label{sec:verify_ppe}
\label{sec:dc_iteration_accuracy}

2 次元 Poisson 問題（参照解 $p^* = \sin\pi x\,\sin\pi y$，Dirichlet BC）に対して
DC 反復回数 $k = 1, 2, 3, 5, 10$ の $L^\infty$ 誤差を $N = 8$--$128$ で計測する．

\begin{table}[htbp]
\centering
\caption{DC 反復回数 $k$ と $L^\infty$ 誤差（Dirichlet BC）}
\label{tab:dc_k_accuracy}
\renewcommand{\arraystretch}{1.2}
\footnotesize
\begin{tabular}{rr ccccc}
\toprule
$N$ & $h$ & $k=1$ & $k=2$ & $k=3$ & $k=5$ & $k=10$ \\
\midrule
  8 & $1/8$   & $1.30 \times 10^{-2}$ & $2.71 \times 10^{-4}$ & $2.08 \times 10^{-4}$ & $2.06 \times 10^{-4}$ & $1.98 \times 10^{-4}$ \\
 16 & $1/16$  & $3.22 \times 10^{-3}$ & $1.13 \times 10^{-5}$ & $1.91 \times 10^{-6}$ & $1.90 \times 10^{-6}$ & $1.81 \times 10^{-6}$ \\
 32 & $1/32$  & $8.04 \times 10^{-4}$ & $6.54 \times 10^{-7}$ & $1.60 \times 10^{-8}$ & $1.59 \times 10^{-8}$ & $1.53 \times 10^{-8}$ \\
 64 & $1/64$  & $2.01 \times 10^{-4}$ & $4.04 \times 10^{-8}$ & $1.29 \times 10^{-10}$ & $1.28 \times 10^{-10}$ & $1.23 \times 10^{-10}$ \\
128 & $1/128$ & $5.02 \times 10^{-5}$ & $2.52 \times 10^{-9}$ & $2.07 \times 10^{-12}$ & $1.59 \times 10^{-12}$ & $2.05 \times 10^{-12}$ \\
\midrule
\multicolumn{2}{c}{収束次数} & $\approx 2.0$ & $\approx 4.0$ & $\approx 7.0$ & $\approx 7.0$ & $\approx 7.0$ \\
\bottomrule
\end{tabular}
\end{table}

$k = 1$ では基底ソルバ $L_L$（$\Ord{h^2}$ FD Laplacian）の打ち切り誤差が支配し
$\Ord{h^2}$ にとどまる．$k = 2$ で中間的な $\Ord{h^4}$ となり，
$k \geq 3$ で CCD の設計精度が完全に発現して $\Ord{h^7}$ を達成する．
設計次数 $\Ord{h^6}$ を超える超収束は，Dirichlet 条件下で参照解が境界上で零となり
境界スキームの $\Ord{h^5}$ 打ち切り誤差が消失することに起因する．
$k = 3$ と $k = 10$ の絶対誤差差は高々 2 倍であり，
全 PPE 求解には $k = 3$ を採用する．
変密度条件下の DC 緩和パラメタ $\omega$ 依存性は
付録~\ref{app:dc_convergence_theory} の解析および §~\ref{sec:U6_split_ppe_dc_hfe} の
分相 PPE で評価する．
\label{sec:verify_dc_spectral}%  backward compat (DC ω 緩和とスペクトル収束)

\noindent\textbf{結論：}$k \geq 3$ で $\Ord{h^{7.0}}$（Dirichlet 超収束）．\textbf{合格}．

% ───────────────
\paragraph{(b) PPE BC：Option III/IV ゲージ + Neumann consistency}
\label{sec:verify_ppe_neumann}

全周 Neumann BC + 1 点ゲージ固定での DC $k = 3$ の格子収束性を，
参照解 $p^* = \cos\pi x\,\cos\pi y$（境界上で自然に $\partial p/\partial n = 0$）で検証する．

\begin{table}[htbp]
\centering
\caption{PPE Neumann BC + ゲージ固定 $p_{0,0} = 1$ の格子収束}
\label{tab:ppe_neumann}
\renewcommand{\arraystretch}{1.2}
\footnotesize
\begin{tabular}{rr cc}
\toprule
$N$ & $h$ & $L^\infty$ 誤差 & 次数 \\
\midrule
  8 & $1/8$   & $1.98 \times 10^{-2}$ & ---   \\
 16 & $1/16$  & $8.59 \times 10^{-4}$ & $4.5$ \\
 32 & $1/32$  & $2.87 \times 10^{-5}$ & $4.9$ \\
 64 & $1/64$  & $9.09 \times 10^{-7}$ & $5.0$ \\
128 & $1/128$ & $2.85 \times 10^{-8}$ & $5.0$ \\
\bottomrule
\end{tabular}
\end{table}

$N \geq 32$ で $\Ord{h^{5.0}}$ に安定する．
これは CCD 内部スキーム（$\Ord{h^6}$）ではなく
境界スキーム（$\Ord{h^5}$）が全体精度のボトルネックであることを意味する．
二相流 PPE は Neumann 条件で運用するため，
達成可能な空間精度は $\Ord{h^5}$ が実質的な上界となる．

\noindent\textbf{結論：}$\Ord{h^{5.0}}$（境界スキーム律速）．\textbf{合格}．

\medskip
\noindent\textbf{U2 補足：参考データ（旧解析の保持）}\\
\noindent\label{sec:verify_ppe_varrho}\label{sec:varrho_ppe_limitation}%
変密度 PPE の参考データ：滑らかな密度場 $\rho(x,y) = 1 + A\sin\pi x\cos\pi y$
（$\rho_{\max}/\rho_{\min} \approx 999$）に対して DC $k=3$ は
$\Ord{h^{5.8}}$ を維持する（$N = 128$ で $L^\infty = 8.5 \times 10^{-11}$）．
界面型 smoothed Heaviside 密度では高解像度側で漸近格子収束が失われる
本質的限界があり，これが分相 PPE 解法（U6）を必要とする根拠である．\\
\noindent\label{sec:verify_pressure_filter}%
圧力場フィルタ適用禁止の定量検証は U9（§~\ref{sec:U9_dccd_pressure_prohibition}）に
配置した（DCCD-on-pressure 禁止則の否定テスト）．\\
\noindent\label{sec:verify_ppe_condition}%
PPE 条件数のスケーリング解析：等密度で $\kappa \propto N^{3.0}$，
$\rho_l/\rho_g = 1000$, $N = 64$ で $\kappa \approx 2.2 \times 10^7$．
分相 PPE による条件数低減の必要性を裏付ける（解析議論は §~\ref{sec:ppe_bc_stability}）．\\
\noindent\label{sec:verify_ccd_spectral_radius}%
CCD $D^{(2)}$ スペクトル半径：$\rho h^2 \approx 9.6$（FD 比 2.4 倍）．
明示的粘性 CFL を $\nu\Delta t/h^2 < 1/(2 \times 9.6) \approx 0.052$ に制約するが，
CN 陰的処理（U8-c）採用によりこの制約は回避される．
